\section{Experimental Setup}
\label{sec:setup}

\parabf{Implementation.} 
While \tech is general across different software agent scaffolds, we implement \tech on top of the popular \minisweagent~\citep{yang2024sweagent} framework. 
As such, we retain the hyperparameters used in \minisweagent by default (i.e., maximum step limit of 250 and maximum cost of \$3 per issue).
We choose \minisweagent to build on since it is not only simplistic (with $\sim$100 lines of code and only accessing bash commands) but also widely used.
Unless otherwise stated, we use \claudesonnetfourfive (\CodeIn{claude-sonnet-4-5-20250929})~\citep{sonnet4.5} in our experiments and sample one patch per issue.


\parabf{Datasets.}
We evaluate \tech on the popular \swebench Verified benchmark~\citep{sbv} containing 500 software development problems where the goal is to successfully modify the repository given a problem description. 
\swebench Verified is validated by human developers to ensure each problem description has sufficient amount of information to solve the issue.
Additionally, we also evaluate on the recent \swebenchpro~\citep{deng2025swe} benchmark, containing 731 publicly available problems, aimed to capture realistic, complex, and enterprise-level problems. 
Compared with \swebench Verified, \swebenchpro contains more difficult problems across multiple repositories and programming languages.

\parabf{Baselines.}
We compare \tech against representative state-of-the-art agentic approaches. 
For \swebench Verified, we compare against \minisweagent as it is one of the most widely-used open-source agentic solutions on \swebench tasks with top leaderboard performance.
Additionally, it is also a straightforward comparison as we directly build on top of the \minisweagent framework.
Furthermore, we compare with prior self-evolving agents: Self-Improving Coding Agent (SICA)~\citep{robeyns2025sica}, Darwin-Gödel Machine (DGM)~\citep{zhang2025darwin}, and Huxley-Gödel Machine (HGM)~\citep{wang2025huxley} on a subset of \swebench Verified problems. 
This subset of 60 \swebench Verified problems has been used by prior work~\citep{wang2025huxley} to specifically evaluate all three self-improving agent baselines.
For \swebenchpro, we compare against \sweagent~\citep{yang2024sweagent} which is the top-performing approach on the \swebenchpro leaderboard.
For each baseline, we directly reuse their experimental results and report their performance, cost, as well as the backend \llm used whenever possible.
